\chapter{热传导}
\section{热传导方程的导出}
目的求出(x,y,z)位置的温度t
\begin{definition}[傅里叶定律]
根据传热学中的傅里叶定律：物体在无穷小时间段 $\mathrm{d}t$ 内，
沿法线方向 $n$ 流过一个无穷小面积 $\mathrm{d}S$ 的热量 $\mathrm{d}Q$ 与
物体温度沿曲面 $\mathrm{d}S$ 法线方向的方向导数 $\dfrac{\partial u}{\partial n}$ 成正比，即
\[
  \mathrm{d}Q=-k(x,y,z)\,\frac{\partial u}{\partial n}\,\mathrm{d}S\,\mathrm{d}t,
\]
其中 $k(x,y,z)$ 称为物体在点 $(x,y,z)$ 处的热传导系数，应取正值。
式中的负号是因为热量总是由高温处流向低温处，因此 $\mathrm{d}Q$ 与
$\dfrac{\partial u}{\partial n}$ 异号。
\end{definition}
\begin{note}
     在时间 $\mathrm{d}t$ 内，通过小面积 $\mathrm{d}S$ 的热量 $\mathrm{d}Q$
与该处沿法线方向的温度梯度成正比，比例系数是热导率 $k(x,y,z)$，
并且热量总是从高温流向低温（所以有负号）。
\end{note}
在物体 $G$ 内任取一闭曲面 $\Gamma$，它所包围的区域记为 $\Omega$。由式(1.1)可知，
从时刻 $t_1$ 到 $t_2$ 流进此闭曲面的全部热量为
\[
  Q=\int_{t_1}^{t_2}\left[\iint_{\Gamma} k(x,y,z)\,\frac{\partial u}{\partial n}\,\mathrm{d}S\right]\mathrm{d}t,
\tag{1.2}
\]
这里 $\dfrac{\partial u}{\partial n}$ 表示 $u$ 沿 $\Gamma$ 上单位外法线方向 $n$ 的方向导数。
\begin{note}
    这里面两个积分号是因为对面积积分
\end{note}
\begin{theorem}[能量守恒关系]
设 $G$ 为物体所占区域，在 $G$ 内任取一闭曲面 $\Gamma$，其所包围的区域记为 $\Omega$。
在时间区间 $(t_1,t_2)$ 内，物体温度由 $u(x,y,z,t_1)$ 变为 $u(x,y,z,t_2)$，则该体积
$\Omega$ 内应吸收的热量为
\[
  \iiint_{\Omega} c(x,y,z)\,\rho(x,y,z)\,
  \bigl[u(x,y,z,t_2)-u(x,y,z,t_1)\bigr]\,\mathrm{d}x\,\mathrm{d}y\,\mathrm{d}z,
\]
其中 $c$ 为比热容，$\rho$ 为密度。由能量守恒可得
\[
  \int_{t_1}^{t_2}\iint_{\Gamma} k(x,y,z)\,\frac{\partial u}{\partial n}\,\mathrm{d}S\,\mathrm{d}t
  = \iiint_{\Omega} c(x,y,z)\,\rho(x,y,z)\,
  \bigl[u(x,y,z,t_2)-u(x,y,z,t_1)\bigr]\mathrm{d}x\,\mathrm{d}y\,\mathrm{d}z.
  \tag{1.3}
\]
\end{theorem}
\begin{theorem}[格林公式的应用]
假设函数 $u$ 关于变量 $x,y,z$ 具有二阶连续偏导数，关于 $t$ 具有一阶连续偏导数。
利用格林公式，可以把式 \textup{(1.3)} 化为
\[
\int_{t_1}^{t_2}\iiint_{\Omega}
\left[
\frac{\partial}{\partial x}\!\left(k\frac{\partial u}{\partial x}\right)
+\frac{\partial}{\partial y}\!\left(k\frac{\partial u}{\partial y}\right)
+\frac{\partial}{\partial z}\!\left(k\frac{\partial u}{\partial z}\right)
\right]\,
\mathrm d x\,\mathrm d y\,\mathrm d z\,\mathrm d t
=
\iiint_{\Omega}
c\rho\left(\int_{t_1}^{t_2}\frac{\partial u}{\partial t}\,\mathrm d t\right)
\mathrm d x\,\mathrm d y\,\mathrm d z .
\]
交换积分次序，就得到
\[
\int_{t_1}^{t_2}\iiint_{\Omega}
\left[
c\rho\frac{\partial u}{\partial t}
-\frac{\partial}{\partial x}\!\left(k\frac{\partial u}{\partial x}\right)
-\frac{\partial}{\partial y}\!\left(k\frac{\partial u}{\partial y}\right)
-\frac{\partial}{\partial z}\!\left(k\frac{\partial u}{\partial z}\right)
\right]\,
\mathrm d x\,\mathrm d y\,\mathrm d z\,\mathrm d t
=0. \tag{1.4}
\]
\end{theorem}
\begin{definition}[热传导方程]
设 $u(x,y,z,t)$ 表示各向同性固体中点 $(x,y,z)$ 在时刻 $t$ 的温度，
$k(x,y,z)$ 为热导率，$c(x,y,z)$ 为比热，$\rho(x,y,z)$ 为密度。
若 $u$ 满足偏微分方程
\begin{equation}\label{eq:nonuniform-heat}
c\rho\,\frac{\partial u}{\partial t}
=\frac{\partial}{\partial x}\!\left(k\frac{\partial u}{\partial x}\right)
+\frac{\partial}{\partial y}\!\left(k\frac{\partial u}{\partial y}\right)
+\frac{\partial}{\partial z}\!\left(k\frac{\partial u}{\partial z}\right),
\end{equation}
则称 \eqref{eq:nonuniform-heat} 为\emph{非均匀各向同性体}的热传导方程（简称热方程）。

如果物体是\emph{均匀}的，此时 $k,c,\rho$ 均为常数，记
\[
\frac{k}{c\rho}=a^{2},
\]
则上式化为
\begin{equation}\label{eq:uniform-heat}
\frac{\partial u}{\partial t}
=a^{2}\left(
\frac{\partial^{2}u}{\partial x^{2}}
+\frac{\partial^{2}u}{\partial y^{2}}
+\frac{\partial^{2}u}{\partial z^{2}}
\right),
\end{equation}
称 \eqref{eq:uniform-heat} 为\emph{均匀各向同性体}的（经典）热传导方程。
\end{definition}
若所考察的物体内有热源（例如物体中通有电流，或有化学反应等情况），则在热传导方程的推导中还需考虑热源的影响。
若设在单位时间内单位体积中所产生的热量为 $F(x,y,z,t)$，则在考虑热平衡时，(1.3) 式左端应再加上一项
\[
\int_{t_1}^{t_2}\iiint_{\Omega} F(x,y,z,t)\,dx\,dy\,dz\,dt.
\]

于是，相应于 (1.6) 式的热传导方程应改为
\begin{equation}\label{eq:heat-nonhomogeneous}
\frac{\partial u}{\partial t}
= a^{2}\left(
\frac{\partial^{2}u}{\partial x^{2}}
+\frac{\partial^{2}u}{\partial y^{2}}
+\frac{\partial^{2}u}{\partial z^{2}}
\right)+f(x,y,z,t), \tag{1.7}
\end{equation}
其中
\begin{equation}\label{eq:f-source}
f(x,y,z,t)=\frac{F(x,y,z,t)}{\rho c}. \tag{1.8}
\end{equation}

(1.6) 式称为齐次热传导方程，而 (1.7) 式称为非齐次热传导方程。
\subsection{定解问题的提法}
\begin{note}
   知道了物体在 \textcolor{purple}{边界上的温度状况}（或热交换状况）
和物体在 \textcolor{purple}{初始时刻的温度}，
就可以确定物体在以后时刻的温度。
因此热传导方程最自然的一个定解问题就是在已给的
\textcolor{purple}{初始条件} 与 \textcolor{purple}{边界条件} 下求问题的解。
\end{note}
\begin{definition}[初始条件]
热传导问题的初始条件提法为
\begin{equation}\label{eq:init}
u(x,y,z,0)=\varphi(x,y,z),
\end{equation}
其中 $\varphi(x,y,z)$ 为已知函数，表示物体在 $t=0$ 时的温度分布。
\end{definition}
\begin{theorem}[热传导方程的三种典型边界条件]
设 $\Omega$ 为物体所占区域，$\Gamma=\partial\Omega$ 为其边界，
$0<t<T$。热传导方程的常见边界条件可分为三类：

\begin{itemize}
  \item[(1)] \textbf{第一类（Dirichlet 边界条件）——给定温度}
  \[
    u(x,y,z,t)\big|_{(x,y,z)\in\Gamma}=g(x,y,z,t).
  \]
  在边界上直接规定温度 $u$ 的数值，温度本身被“强制”固定；
  主要适用于边界与温度已知的巨大热源或恒温器接触的情形。

  \item[(2)] \textbf{第二类（Neumann 边界条件）——给定热流/法向导数}
  \[
    \frac{\partial u}{\partial n}\Big|_{(x,y,z)\in\Gamma}=g(x,y,z,t).
  \]
  在边界上规定温度沿外法线方向的导数（等价于给定单位面积上的热流密度），
  此时影响最大的是法向导数这一项；典型情形包括：绝热边界
 （$g\equiv 0$）或热流强度预先已知的边界。

  \item[(3)] \textbf{第三类（Robin 边界条件）——与环境换热}
  \[
    \Big(\frac{\partial u}{\partial n}+\sigma u\Big)\Big|_{(x,y,z)\in\Gamma}
      = g(x,y,z,t), \qquad \sigma>0.
  \]
  边界上的热流既由内部传导决定，又由与外界介质的换热决定，
  其中换热系数 $\sigma$ 控制哪一部分影响更大：
  当 $\sigma$ 很大时，温度更接近被外界环境“拉”到某一值（类似第一类）；
  当 $\sigma$ 很小时，边界更接近给定热流或绝热的情况（类似第二类）。
\end{itemize}
\end{theorem}
\begin{definition}[热传导方程的柯西问题]
若所考察的物体体积很大，而我们只关心在较短时间和较小区域内的温度变化，
边界条件的影响可以忽略，此时可把物体视为充满整个空间。
对应的热传导方程定解问题称为\emph{柯西问题}，其初始条件为
\begin{equation}\label{eq:heat-cauchy}
u(x,y,z,0)=\varphi(x,y,z),\qquad -\infty<x,y,z<\infty. \tag{1.14}
\end{equation}
这里 $\varphi(x,y,z)$ 为已知函数，表示在 $t=0$ 时整个空间中的温度分布。
\end{definition}
\begin{note}
    热传导方程的柯西问题：把物体看作充满整个空间，只给出 $t=0$ 时全空间的温度分布作为初始条件，在无边界影响下求以后的温度。
\end{note}
\begin{theorem}[ 一维vs二维]
在适当情形下，方程中描述空间坐标的独立变量个数可以减少，典型情况如下：
\begin{itemize}
  \item[(1)] \textbf{一维热传导方程}\\
  若物体为均匀细杆，其侧面绝热（不与外界发生热交换），并且同一截面上的温度分布相同，
  则温度仅依赖于坐标 $x$ 与时间 $t$，满足
  \begin{equation}\label{eq:1d-heat}
    \frac{\partial u}{\partial t} = a^{2}\frac{\partial^{2}u}{\partial x^{2}}. \tag{1.15}
  \end{equation}

  \item[(2)] \textbf{二维热传导方程}\\
  若考虑薄片的热传导，且薄片的侧面绝热，则温度仅依赖于平面坐标 $x,y$ 与时间 $t$，
  满足
  \begin{equation}\label{eq:2d-heat}
    \frac{\partial u}{\partial t}
    = a^{2}\left(
        \frac{\partial^{2}u}{\partial x^{2}}
       +\frac{\partial^{2}u}{\partial y^{2}}
      \right). \tag{1.16}
  \end{equation}
\end{itemize}
\end{theorem}
\begin{note}
    一维是只考虑温度沿一条线（一个空间坐标）的变化，而二维是要同时考虑温度在一个平面上（两个空间坐标）的扩散与变化
\end{note}
\subsection{扩散问题}
\begin{itemize}
  \item 在分子扩散问题中，建立描述扩散过程的数学模型。
  \item 研究扩散物质的浓度这一主要量（单位体积中所含扩散物质的质量）。
  \item 利用与热传导过程相似的物理规律，类比热传导方程的推导，写出扩散方程。
\end{itemize}
\begin{definition}[傅里叶定律与热量守恒定律]
在推导热传导方程的过程中，起基本作用的是下列两条物理定律：

\begin{itemize}
  \item 傅里叶定律：
  \[
    \mathrm dQ = -k(x,y,z)\,\frac{\partial u}{\partial n}\,\mathrm dS\,\mathrm dt,
  \]
  其中 $u$ 为温度，$k(x,y,z)$ 为热导率，$\dfrac{\partial u}{\partial n}$ 为沿外法线方向的导数。

  \item 热量守恒定律：
  \[
    \int_{t_1}^{t_2}\iint_{\Gamma} k\,\frac{\partial u}{\partial n}\,\mathrm dS\,\mathrm dt
    =\iiint_{\Omega} c\rho\,[u(x,y,z,t_2)-u(x,y,z,t_1)]\,\mathrm d x\,\mathrm d y\,\mathrm d z,
  \]
  其中 $c$ 为比热，$\rho$ 为密度，$\Omega$ 为物体所占区域，$\Gamma=\partial\Omega$ 为其边界。
\end{itemize}
\end{definition}
\begin{itemize}
  \item 傅里叶定律：描述热量在微观上如何沿温度梯度传递——温度沿法线方向变化越快、热导率越大，单位时间通过单位面积的热流就越大，并且热量总是从高温流向低温。
  \item 热量守恒定律：说明在封闭区域内，某时间段内内部总热量的增量，正好等于这段时间内通过边界流入（或流出）的全部热量。
\end{itemize}
\begin{definition}[反应扩散方程]
若在区域 $\Omega$ 中存在产生扩散物质的源（例如化学反应），
并且在单位时间内产生密度为 $f(x,y,z,t)$ 的扩散物质，则扩散过程满足的方程为
\begin{equation}\label{eq:reaction-diffusion}
\frac{\partial N}{\partial t}
=\frac{\partial}{\partial x}\!\left(D\frac{\partial N}{\partial x}\right)
+\frac{\partial}{\partial y}\!\left(D\frac{\partial N}{\partial y}\right)
+\frac{\partial}{\partial z}\!\left(D\frac{\partial N}{\partial z}\right)
+f(x,y,z,t), \tag{1.20}
\end{equation}
称为\emph{反应扩散方程}，其中 $N(x,y,z,t)$ 为扩散物质的浓度，$D$ 为扩散系数。
\end{definition}
% 扩散方程的推导（分步）

% 扩散方程的推导（分步）

\begin{description}

  \item[步骤1：扩散定律（Fick 定律）——建立通过边界的小量关系]
  以 $N(x,y,z,t)$ 表示扩散物质的浓度，$dm$ 表示在无穷小时间段 $dt$
  内沿曲面 $S$ 的法向 $n$ 经过面积元 $dS$ 的扩散物质量，$D(x,y,z)$ 为扩散系数，
  则有扩散定律
  \begin{equation}\label{eq:fick}
    dm = -D(x,y,z)\,\frac{\partial N}{\partial n}\,dS\,dt. \tag{1.17}
  \end{equation}

  \item[步骤2：质量守恒——把边界通量与总体质量变化联系起来]
  与热传导过程类似，在任意有界区域 $\Omega$ 中关于扩散物质的质量守恒定律给出
  \begin{equation}\label{eq:mass-cons}
    \int_{t_1}^{t_2}\iint_{\Gamma}
      D\,\frac{\partial N}{\partial n}\,dS\,dt
    =\iiint_{\Omega}\bigl[N(x,y,z,t_2)-N(x,y,z,t_1)\bigr]\,
      dx\,dy\,dz, \tag{1.18}
  \end{equation}
  其中 $\Gamma=\partial\Omega$。

  \item[步骤3：类比热传导——从积分守恒律得到微分形式的扩散方程]
  将 \eqref{eq:fick}、\eqref{eq:mass-cons} 与热传导中的
  傅里叶定律和热量守恒式对比，可知二者在形式上完全类似；
  只需用 $m,N,D$ 分别对应热传导中的 $Q,u,k$，并注意到此处的 $c\rho$
  相当于常数 $1$，即可写出扩散方程
  \begin{equation}\label{eq:diffusion}
    \frac{\partial N}{\partial t}
    =\frac{\partial}{\partial x}\!\left(D\frac{\partial N}{\partial x}\right)
    +\frac{\partial}{\partial y}\!\left(D\frac{\partial N}{\partial y}\right)
    +\frac{\partial}{\partial z}\!\left(D\frac{\partial N}{\partial z}\right).
    \tag{1.19}
  \end{equation}

  \item[步骤4：常系数情形——得到与热方程同型的扩散方程]
  若扩散系数 $D$ 为常数，记 $D=a^{2}$，则 \eqref{eq:diffusion} 化为
  \[
    \frac{\partial N}{\partial t}
    =a^{2}\left(
      \frac{\partial^{2}N}{\partial x^{2}}
      +\frac{\partial^{2}N}{\partial y^{2}}
      +\frac{\partial^{2}N}{\partial z^{2}}
    \right),
  \]
  其形式与均匀各向同性体的热传导方程完全相同。
\end{description}
\section{初边值的分离变量问题}

考虑典型问题
\[
\begin{cases}
u_t = a^2 u_{xx}, & t>0,\ 0<x<l,\\[2pt]
u(x,0)=\varphi(x), & 0<x<l,\\[2pt]
u(0,t)=0,\quad u(l,t)=0, & t>0.
\end{cases}
\]

\textbf{一、什么时候用分离变量法？}
\begin{itemize}
  \item 线性、齐次的偏微分方程（热方程、波动方程、拉普拉斯方程等）。
  \item 区域形状简单（区间、矩形、长方体等），边界条件只依赖空间，不随时间变。
  \item 边界条件最好是齐次的（如 $u=0$ 或 $u_n=0$）；若不齐次，通常先减去一个稳态解把边界条件变成齐次，再用分离变量法。
\end{itemize}

\textbf{二、怎么用分离变量法？（一维热方程的标准步骤）}
\begin{enumerate}
  \item \emph{假设可分离解}：设
  \[
    u(x,t)=X(x)T(t),
  \]
  代入方程 $u_t=a^2u_{xx}$ 得
  \[
    X T' = a^2 X'' T \quad\Rightarrow\quad
    \frac{T'}{a^2 T} = \frac{X''}{X} = -\lambda.
  \]

  \item \emph{得到两个常微分方程}：
  \[
    T' + a^2\lambda T=0,\qquad X'' + \lambda X =0.
  \]

  \item \emph{用边界条件求特征值与特征函数}：

  边界条件 $u(0,t)=u(l,t)=0$ 变成
  \[
    X(0)=0,\quad X(l)=0.
  \]
  解 Sturm–Liouville 问题
  \[
    X''+\lambda X=0,\quad X(0)=X(l)=0
  \]
  得到
  \[
    \lambda_k=\frac{k^2\pi^2}{l^2},\qquad
    X_k(x)=\sin\frac{k\pi x}{l},\qquad k=1,2,\dots
  \]

  \item \emph{解时间方程并线性叠加}：

  每个 $\lambda_k$ 对应一个
  \[
    T_k(t)=e^{-a^2\lambda_k t}=e^{-a^2 k^2\pi^2 t/l^2}.
  \]
  线性方程可叠加，所以
  \[
    u(x,t)=\sum_{k=1}^\infty A_k e^{-a^2 k^2\pi^2 t/l^2}\sin\frac{k\pi x}{l}.
  \]

  \item \emph{用初始条件求系数 $A_k$}：

  由 $u(x,0)=\varphi(x)$ 得
  \[
    \varphi(x)=\sum_{k=1}^\infty A_k\sin\frac{k\pi x}{l},
  \]
  即是 $\varphi$ 在 $(0,l)$ 上的正弦傅里叶级数展开，系数为
  \[
    A_k=\frac{2}{l}\int_0^l \varphi(x)\sin\frac{k\pi x}{l}\,dx.
  \]
\end{enumerate}

\textbf{三、为什么这样做是合理的？（做题用到的大致理由）}
\begin{itemize}
  \item 热方程是线性齐次方程，满足叠加原理：单个分离解可以线性组合成一般解。
  \item 带边界条件的空间方程 $X''+\lambda X=0$ 是特征值问题，其特征函数
        $\{\sin(k\pi x/l)\}$ 在 $(0,l)$ 上构成一组完备正交基，可以展开任意“够好”的函数
        $\varphi(x)$。
  \item 时间因子 $e^{-a^2\lambda_k t}$ 体现了各个模式的衰减，不影响求系数时的正交性，
        所以先展开 $\varphi$、再附上时间因子即可得到满足 PDE+边界+初始条件的解。
\end{itemize}

\textbf{最终记住的模板：}
\[
\boxed{\;
u(x,t)=\sum_{k=1}^\infty
\left(\frac{2}{l}\int_0^l \varphi(x)\sin\frac{k\pi x}{l}\,dx\right)
e^{-a^2 k^2\pi^2 t/l^2}\sin\frac{k\pi x}{l}
\;}
\]
只要题目是同样的边界条件 $u(0,t)=u(l,t)=0$，换不同的 $\varphi(x)$，
直接套这个模板计算 $A_k$ 就可以做作业题。
\begin{homework}[2]
用分离变量法求解热传导方程的初边值问题：
\[
\begin{cases}
\dfrac{\partial u}{\partial t}
= \dfrac{\partial^2 u}{\partial x^2}, & (t>0,\,0<x<1),\\[6pt]
u(x,0)=
\begin{cases}
x, & 0<x\le\dfrac{1}{2},\\[4pt]
1-x, & \dfrac{1}{2}<x<1,
\end{cases}\\[10pt]
u(0,t)=u(1,t)=0, & (t>0).
\end{cases}
\]
\end{homework}
\begin{solution}
 % 题 2
\[
u_t=u_{xx},\qquad u(0,t)=u(1,t)=0,\qquad 
u(x,0)=
\begin{cases}
x,&0<x\le \frac12,\\[2pt]
1-x,&\frac12<x<1
\end{cases}
\]

\[
u(x,t)=X(x)T(t)
\]

\[
\frac{T'}{T}=\frac{X''}{X}=-\lambda
\]

\[
T'+\lambda T=0,\qquad T(t)=C\,e^{-\lambda t}
\]

\[
X''+\lambda X=0,\qquad X(0)=X(1)=0
\]

\[
\lambda=(n\pi)^2,\qquad X_n(x)=\sin(n\pi x),\qquad n=1,2,\ldots
\]

\[
u(x,t)=\sum_{n=1}^{\infty}b_n\,e^{-(n\pi)^2 t}\,\sin(n\pi x)
\]

\[
b_n=2\int_{0}^{1}u(x,0)\sin(n\pi x)\,dx
=2\!\left(\int_{0}^{1/2}x\sin(n\pi x)\,dx+\int_{1/2}^{1}(1-x)\sin(n\pi x)\,dx\right)
\]

\[
\int x\sin(\alpha x)\,dx=-\frac{x\cos(\alpha x)}{\alpha}+\frac{\sin(\alpha x)}{\alpha^2}
\]

\[
b_n=\frac{2}{n\pi}\Bigl[-x\cos(n\pi x)\Bigr]_{0}^{1/2}
+\frac{2}{(n\pi)^2}\Bigl[\sin(n\pi x)\Bigr]_{0}^{1/2}
+\frac{2}{n\pi}\Bigl[-(1-x)\cos(n\pi x)\Bigr]_{1/2}^{1}
-\frac{2}{(n\pi)^2}\Bigl[\sin(n\pi x)\Bigr]_{1/2}^{1}
\]

\[
b_n=\frac{4}{(n\pi)^2}\,\sin\!\Bigl(\frac{n\pi}{2}\Bigr)
\]

\[
u(x,t)=\sum_{n=1}^{\infty}\frac{4}{\pi^2 n^2}\,\sin\!\Bigl(\frac{n\pi}{2}\Bigr)\,e^{-(n\pi)^2 t}\,\sin(n\pi x)
\]   
\end{solution}
\begin{dxtips}
 \begin{itemize}
  \item \textbf{分离变量原理}\\
  线性齐次偏微分方程在齐次边界条件下可以尝试形如
  $u(x,t)=X(x)T(t)$ 的乘积解，从而将 PDE 化为两个常微分方程。

  \item \textbf{自伴随算子的特征值问题}\\
  空间算子 $L[X]=X''$ 加齐次狄利克雷边界条件 $X(0)=X(1)=0$ 构成
  Sturm–Liouville 系统，其特征值与特征函数为
  \[
    \lambda_n=(n\pi)^2,\qquad X_n(x)=\sin(n\pi x).
  \]

  \item \textbf{正交基展开原理}\\
  $\{\sin(n\pi x)\}_{n\ge1}$ 作为 Sturm–Liouville 系统的特征函数，
  在 $L^2(0,1)$ 中构成一组完备正交基，可用于展开任意初始函数
  $f\in L^2(0,1)$。

  \item \textbf{傅里叶正弦级数的系数公式}\\
  基于正交性
  \[
    \int_0^1 \sin(n\pi x)\sin(m\pi x)\,dx = \frac12\delta_{mn},
  \]
  得到初始条件展开系数
  \[
    b_n = 2\int_0^1 f(x)\sin(n\pi x)\,dx.
  \]

  \item \textbf{线性叠加与唯一性原理}\\
  对线性齐次 PDE，所有特征解的级数叠加仍是解；
  热方程满足适定性与唯一性，故由级数展开得到的解即为唯一解。

\end{itemize}
\end{dxtips}
\section{柯西问题}
\subsection{傅立叶}
% 傅里叶变换的由来（作业用分步梳理）

\begin{description}

  \item[步骤1：从傅里叶级数出发——在有限区间上展开函数]
  设 $f(x)$ 定义在 $(-\infty,\infty)$ 上，在区间 $[-l,l]$ 上有一阶连续导数，
  则在 $(-l,l)$ 中可以展开为傅里叶级数
  \begin{equation}\label{eq:FS}
    f(x)=\frac{a_0}{2}+\sum_{n=1}^{\infty}
        \Bigl(a_n\cos\frac{n\pi}{l}x
             +b_n\sin\frac{n\pi}{l}x\Bigr), \tag{3.1}
  \end{equation}
  其中系数
  \begin{equation}\label{eq:FS-coef}
    a_n=\frac{1}{l}\int_{-l}^{l}f(\xi)\cos\frac{n\pi}{l}\xi\,d\xi,\qquad
    b_n=\frac{1}{l}\int_{-l}^{l}f(\xi)\sin\frac{n\pi}{l}\xi\,d\xi. \tag{3.2}
  \end{equation}
  \emph{目的：}把 $f(x)$ 在有限区间上写成正弦、余弦的线性组合，为后来过渡到“连续频率”做准备。

  \item[步骤2：把级数改写成“和的积分”形式——为极限 (和$\to$积分) 做准备]
  将 \eqref{eq:FS-coef} 代入 \eqref{eq:FS}，可整理为
  \begin{equation}\label{eq:pre-integral}
    f(x)=\frac{1}{2l}\int_{-l}^{l}f(\xi)\,d\xi
         +\sum_{n=1}^{\infty}\frac{1}{l}\int_{-l}^{l}
           f(\xi)\cos\frac{n\pi}{l}(x-\xi)\,d\xi.
  \end{equation}
  \emph{目的：}把“傅里叶级数的和”写成“权重为 $\frac{1}{l}$ 的积分的和”，
  以后就可以把这类和看成某个积分的黎曼和。

  \item[步骤3：令区间扩展到整个实轴——得到傅里叶积分表示]
  若 $f(x)$ 在 $(-\infty,\infty)$ 上绝对可积，令 $l\to\infty$，
  记
  \[
    \lambda_n=\frac{n\pi}{l},\qquad
    \Delta\lambda=\lambda_{n+1}-\lambda_n=\frac{\pi}{l},
  \]
  则 \eqref{eq:pre-integral} 在极限意义下变为
  \begin{equation}\label{eq:Fourier-integral}
    f(x)=\frac{1}{\pi}\int_{0}^{\infty} d\lambda
         \int_{-\infty}^{\infty} f(\xi)\cos\lambda(x-\xi)\,d\xi. \tag{3.3}
  \end{equation}
  \emph{目的：}把离散频率 $n\pi/l$ 的级数极限，变成连续频率 $\lambda$ 的积分，
  得到 $f(x)$ 的傅里叶积分表达式（实数余弦形式）。

  \item[步骤4：改写成复指数形式——引出变换核 $e^{i\lambda x}$]
  由于 $\cos\lambda(x-\xi)$ 是 $\lambda$ 的偶函数，$\sin\lambda(x-\xi)$ 是奇函数，
  可以把 \eqref{eq:Fourier-integral} 写成复指数形式
  \begin{equation}\label{eq:complex-form}
  \begin{aligned}
    f(x)
      &=\frac{1}{2\pi}\int_{-\infty}^{\infty} d\lambda
        \int_{-\infty}^{\infty} f(\xi)
           \bigl[\cos\lambda(x-\xi)+i\sin\lambda(x-\xi)\bigr]\,d\xi \\
      &=\frac{1}{2\pi}\int_{-\infty}^{\infty} d\lambda
        \int_{-\infty}^{\infty} f(\xi)e^{i\lambda(x-\xi)}\,d\xi. \tag{3.4}
  \end{aligned}
  \end{equation}
  \emph{目的：}用 $e^{i\lambda(x-\xi)}$ 统一表示正弦、余弦，方便定义“变换对”。

  \item[步骤5：定义傅里叶变换与逆变换——得到最终的变换公式]
  令
  \begin{equation}\label{eq:FT}
    g(\lambda)=\int_{-\infty}^{\infty}f(\xi)e^{-i\lambda\xi}\,d\xi, \tag{3.5}
  \end{equation}
  则由 \eqref{eq:complex-form} 得到
  \begin{equation}\label{eq:IFT}
    f(x)=\frac{1}{2\pi}\int_{-\infty}^{\infty}g(\lambda)e^{i\lambda x}\,d\lambda. \tag{3.6}
  \end{equation}
  通常称 $g(\lambda)$ 为 $f(x)$ 的\emph{傅里叶变换}，记为 $g=F[f]$；
  称 $f(x)$ 为 $g(\lambda)$ 的\emph{傅里叶逆变换}，记为 $f=F^{-1}[g]$。
  当 $f(x)$ 在 $(-\infty,\infty)$ 上连续可导且绝对可积时，其傅里叶变换存在，
  并且逆变换等于原函数 $f(x)$。
  \emph{目的：}给出真正做题时要用到的“变换对”：
  \[
    F[f](\lambda)=\int_{-\infty}^{\infty}f(x)e^{-i\lambda x}\,dx,\qquad
    F^{-1}[g](x)=\frac{1}{2\pi}\int_{-\infty}^{\infty}g(\lambda)e^{i\lambda x}\,d\lambda.
  \]
\end{description}
\begin{definition}[卷积的定义：用平移后相乘再积分来“混合”两个函数]
设 $f_1,f_2$ 在 $(-\infty,\infty)$ 上绝对可积，定义它们的卷积为
\[
  (f_1*f_2)(x)=\int_{-\infty}^{\infty} f_1(x-t)f_2(t)\,dt,
\]
记作 $f_1*f_2$。若积分存在，则称 $f(x)=f_1*f_2$ 为 $f_1$ 与 $f_2$ 的卷积。
\end{definition}

\begin{theorem}[线性性：先线性组合再变换等于先变换再线性组合]\label{thm:FT-linearity}
若 $f_1,f_2$ 都可以进行傅里叶变换，$\alpha,\beta\in\mathbb{C}$，则
\[
  F[\alpha f_1+\beta f_2](\lambda)
  =\alpha F[f_1](\lambda)+\beta F[f_2](\lambda).
\]
\end{theorem}

\begin{theorem}[卷积定理：卷积在频域里变成简单的乘法]\label{thm:FT-convolution}
若 $f_1,f_2$ 及它们的卷积 $f_1*f_2$ 都可以进行傅里叶变换，则有
\[
  F[f_1*f_2](\lambda)=F[f_1](\lambda)\,F[f_2](\lambda).
\]
\end{theorem}

\begin{theorem}[乘积公式：时域里相乘在频域里变成卷积并多出 $1/(2\pi)$ 因子]\label{thm:FT-product}
若 $f_1,f_2$ 及乘积 $f_1f_2$ 都可以进行傅里叶变换，则
\[
  F[f_1\cdot f_2](\lambda)
  =\frac{1}{2\pi}\,F[f_1](\lambda)*F[f_2](\lambda),
\]
其中右端的 $*$ 表示关于变量 $\lambda$ 的卷积。
\end{theorem}

\begin{theorem}[导数公式：求导在频域中就是乘以 $i\lambda$]\label{thm:FT-derivative}
若 $f,f'$ 都可以进行傅里叶变换，并且当 $|x|\to\infty$ 时 $f(x)\to0$，则
\[
  F[f'](\lambda)=i\lambda\,F[f](\lambda).
\]
\end{theorem}

\begin{theorem}[乘以 $x$ 的公式：时域乘以 $x$ 在频域中对应对变换求导]\label{thm:FT-x-times}
若 $f,\,xf(x)$ 都可以进行傅里叶变换，则
\[
  F[-ix f(x)](\lambda)=\frac{d}{d\lambda}F[f](\lambda).
\]
\end{theorem}
\subsection{极大值原理}
\begin{dxtips}
    看成t为横坐标，x的位置为纵坐标就可以理解为什么叫边了
    
\end{dxtips}
\begin{theorem}[极值原理]\label{thm:max-principle}
设函数 $u(x,t)$ 在矩形
\[
R_T=\{(x,t)\mid \alpha \le x \le \beta,\ 0\le t\le T\}
\]
上连续，并且在矩形内部满足热传导方程
\begin{equation}\label{eq:heat-equation}
  \frac{\partial u}{\partial t}
  = a^{2}\frac{\partial^{2}u}{\partial x^{2}}
  \tag{4.1}
\end{equation}
则 $u$ 在矩形的两个侧边
\[
x=\alpha,\ 0\le t\le T;\qquad x=\beta,\ 0\le t\le T
\]
以及底边
\[
t=0,\ \alpha\le x\le\beta
\]
上取得其最大值和最小值。

换言之，若用 $\Gamma_T$ 表示由上述两侧边及底边所组成的边界曲线（称为抛物边界），则有
\[
\max_{R_T} u(x,t)=\max_{\Gamma_T} u(x,t),\qquad
\min_{R_T} u(x,t)=\min_{\Gamma_T} u(x,t).
\]
\end{theorem}
\begin{itemize}
  \item \textbf{使用条件：}
    \begin{itemize}
      \item 域为矩形 
      \(
      R_T=\{(x,t)\mid \alpha \le x \le \beta,\ 0\le t\le T\}.
      \)
      \item $u$ 在闭矩形 $\overline{R_T}$ 上连续，在内部 $R_T^\circ$ 上足够光滑，使其满足热传导方程 \eqref{eq:heat-equation}。
      \item 抛物边界 $\Gamma_T$ 只包括两侧边 \(x=\alpha,\ x=\beta\) 与底边 \(t=0\)，不包括顶边 \(t=T\)。
    \end{itemize}
  \item \textbf{使用结论：}
    \begin{itemize}
      \item 在整个矩形 $R_T$ 上，$u$ 的最大值、最小值都在抛物边界 $\Gamma_T$ 上取得：
      \(
      \max_{R_T} u = \max_{\Gamma_T} u,\ 
      \min_{R_T} u = \min_{\Gamma_T} u.
      \)
    \end{itemize}
\end{itemize}
\subsection{热核公式}
可以把热核看成一句话：

\[
\text{“单位瞬时点热源，过了时间 }t\text{ 之后在空间中的温度分布。”}
\]

\begin{enumerate}
  \item \textbf{热核 = 点源的“温度指纹”}

  以 $n$ 维热核
  \[
  G_n(x,t)=\frac1{(4\pi a^2 t)^{n/2}}\exp\!\Bigl(-\frac{|x|^2}{4a^2 t}\Bigr)
  \]
  为例：
  \begin{itemize}
    \item 在 $t=0$ 时，把全部热量 $1$ 集中在原点 $x=0$（一个“单位点热源”）；
    \item $G_n(x,t)$ 告诉你：过了时间 $t$，这个点热源在位置 $x$ 处的温度是多少；
    \item 因为 $\displaystyle\int_{\mathbb R^n} G_n(x,t)\,dx = 1$，所以总热量守恒，只是随着时间在空间中扩散开来。
  \end{itemize}

  形象地说：最初是一根“针扎”的温度峰，随后变成越来越平、越来越宽的高斯小山包。

  \item \textbf{卷积公式 = “叠加很多点源”}

  实际初始温度不是一个点，而是一个分布 $\varphi(x)$。物理上可以理解为：
  \begin{itemize}
    \item 每个位置 $y$ 上有一个“点热源”，强度是 $\varphi(y)$；
    \item 每个点热源在时间 $t$ 后对位置 $x$ 的贡献是 $G_n(x-y,t)\,\varphi(y)$；
    \item 把所有点的贡献线性叠加（热方程是线性的），得到
          \[
          u(x,t)=\int_{\mathbb R^n} G_n(x-y,t)\,\varphi(y)\,dy.
          \]
  \end{itemize}
  所以：卷积公式就是“把初始温度看成无数点源，加起来的结果”。

  \item \textbf{记忆小口诀}
  \begin{itemize}
    \item 热核：点热源随时间扩散的“高斯小山”；
    \item 卷积：把初值拆成点热源，利用线性叠加，把所有小山加起来。
  \end{itemize}
\end{enumerate}
\begin{dxtips}
    我们有公式
\[
u(x,t) = \bigl(G_n(\cdot,t)*\varphi\bigr)(x).
\]
可以拆成三部分理解：

\end{dxtips}
\begin{itemize}
  \item \textbf{\(G_n(\cdot,t)\)}：  
  “\(\cdot\)” 只是一个占位符（哑变量），表示把 \(G_n\) 看成关于空间变量的函数。  
  例如也可以写成 \(G_n(y,t)\)，这里只是不指定字母，用 “\(\cdot\)” 占位。

  \item \textbf{\(*\)}（卷积号）：  
  表示两个关于空间变量的函数做卷积。在 \(\mathbb{R}^n\) 上的定义是
  \[
  (f*g)(x) = \int_{\mathbb{R}^n} f(x-y)\,g(y)\,dy.
  \]

  \item \textbf{\(\bigl(G_n(\cdot,t)*\varphi\bigr)(x)\)}：  
  先把 \(G_n(\cdot,t)\) 和 \(\varphi\) 当作两个函数做卷积，得到新的函数
  \(x\mapsto (G_n*\varphi)(x)\)。  
  外面再加上 “\((x)\)” 表示在点 \(x\) 处取值，展开就是
  \[
  \bigl(G_n(\cdot,t)*\varphi\bigr)(x)
  =\int_{\mathbb{R}^n} G_n(x-y,t)\,\varphi(y)\,dy.
  \]
\end{itemize}

所以一句话：
\[
\bigl(G_n(\cdot,t)*\varphi\bigr)(x)
\]
就是“用热核 \(G_n\) 和初值 \(\varphi\) 卷积后，在点 \(x\) 的取值”，也就是上面的积分表达式。 
\begin{dxtips}
    把“满足热传导方程、初值为 \(\varphi\) 的温度场 \(u(x,t)\) 的解”，
转化成了“对初始温度分布 \(\varphi(y)\) 的加权空间积分”——
权函数就是点热源的温度响应 \(G_n(x-y,t)\)。

也就是说，\(u(x,t)\) 被写成
\[
u(x,t)
= \int_{\mathbb{R}^n} G_n(x-y,t)\,\varphi(y)\,dy
= \text{“每个初始点 } y \text{ 上的温度 } \varphi(y)
\text{ 产生一个点热源，通过 } G_n \text{ 传到 } x \text{ 的贡献，}
\text{然后把所有点的贡献积分（叠加）起来”，}
\]
从而把 PDE 的解用一个卷积积分表示出来。
\end{dxtips}